\scanpage{148}
\mspage{6}
\deleted{$\unsure{\zeta(x',u)|\zeta(x',x)}$}\inserted{$x\xrightarrow{u}\alpha(x')$} en prenant $u^*(\zeta(x'))$ (où $\zeta(x')\in(\alpha^*(F))(x')=F(\alpha(x'))$). Le pt essentiel est de montrer que \deleted{\unsure{cet}} les $u^*(\zeta(x'))$ \deleted{\unsure{est}} ne dépendent que de $x$, pas du choix de $(x',u)\in {}_{x\backslash}C'$. Or pour $(x',u)$ variable dans $ {}_{x\backslash}C'$, l'application
\[
\begin{aligned}
\varphi:\ \operatorname{Ob}({}_{x\backslash}C')&\longrightarrow F(x)\\
(x',u)&\longrightarrow u^*(\zeta(x'))
\end{aligned}
\]
est telle que s'il existe une flèche
\[
(x',u)\xrightarrow{u'}(\overline{x'},\bar u)
\qquad\text{où }u':x'\unsure{-|\to}\overline{x'}
\]
\deleted{d'où le diagr} (donc le diagramme
\diagram{\begin{tikzcd}[ampersand replacement=\&]
x\arrow[r,"u"]\arrow[dr,"\bar u"'] \& \alpha(x')\arrow[d,"\alpha(u')"]\\
\&\alpha(\overline{x'})
\end{tikzcd}}
est commutatif), alors $\varphi(x',u)=\varphi(\overline{x'},\bar u)$. Donc $\varphi$ est \unsure{constante} sur les comp. connexes de $ {}_{x\backslash}C'$, donc constante si $ {}_{x\backslash}C'$ est 0-connexe, d'où $\zeta(\unsure{x})$. On prouve alors aussitôt que pour $x$ variable, le
\[
\zeta\overset{\mathrm{def}}{=}(\zeta(x))_{x\in\operatorname{Ob}C}\unsure{\simeq|\in}\varprojlim F,\qquad\deleted{\unsure{x\mapsto\zeta}}.
\]
et que $\alpha^*(\zeta)=\zeta'$, cqfd.

On peut dire que pour $x_0\in\operatorname{Ob}\,C_{\unsure{0}}$ fixé, si on regarde le préfaisceau \uline{constant} sur $ {}_{x_0\backslash}C'$ de valeur $F(x_0)$, \deleted{\illegible{2}} soit $\unsure{\underline{\Phi}_{x_0}|\underline{F}_{x_0}}$, l'application
\[
(x',u)\longmapsto \zeta\unsure{(u')|^{*}(u')|(x')}\quad
\deleted{\operatorname{Ob}({}_{x_0\backslash}C')\longrightarrow F}\ 
{}_{x_0\backslash}C'- F(x_0)
\]
peut être regardé comme une section de ce préfaisceau sur $ {}_{x_0\backslash}C'$, définie par $\zeta'\in\varprojlim_{C'}\alpha^*F$, donc on a
\[
\varprojlim_{C'}\alpha^*F\xrightarrow{c_{x_0}}
\Gamma({}_{x_0\backslash}C'\deleted{\unsure{\zeta}},\underline{\Phi}_{x_0})
\qquad (\simeq F(x_0)\ \text{si }{}_{x_0\backslash}C\ \text{0-conn.})
\]
Si on a $\boxed{x_0\xrightarrow{v}x_1}$ dans $C$, on en déduit
\[
{}_{x_1\backslash}C'\xrightarrow{V^*}{}_{x_0\backslash}C'
\]
(les $ {}_{x\backslash}C'$ sont les fibres d'une \uline{cat. fibrée}\inserted{$ {}_{C\backslash}C'$} sur $C$\ldots) et on a $F(x_1)-F(x_0)$, d'où $\boxed{\underline{\Phi}_{x_1}\longrightarrow V^*\underline{\Phi}_{x_0}}$ pour les faisceaux constants corr. sur $ {}_{x_1\backslash}C'$ resp. $ {}_{x_0\backslash}C'$.
% archive sheet number: 148
